\scanpage{16}
\mspage{14}
$3^\circ)$ Supposons $S=\operatorname{Spec}(\mathbb C)$, prenons $F=T_{\mathrm{Betti}}$ \deleted{\unsure{ou} $T_{\mathrm{Hdg}}$} \deleted{\unsure{foncteur Hdg|Bett / Hdg}.} Alors on \deleted{peut prendre} $K=Q$. D'autre part $F\otimes_Q\mathbb C\simeq T_{\mathrm{Hdg}}\deleted{\unsure{\otimes_Q\mathbb C}}$, donc \deleted{\unsure{on|on a}} $G^{\mathrm{Betti}}\otimes_Q\mathbb C\simeq G^{\mathrm{Hdg}}\deleted{\unsure{\otimes_Q\mathbb C}}$. Donc dans $G^{\mathrm{Betti}}\otimes_Q\mathbb C$ on trouve \deleted{\unsure{un}} $i_1,i_2$
\[
i_1,i_2:\mathrm G_m\longrightarrow G^{\mathrm{Betti}}\otimes_Q\mathbb C
\]
tels que
\[
i_1(\lambda)i_2(\lambda)=\unsure{\eta|\gamma|y}(\lambda).
\]
De plus on a
\[
\overline{i_1(\lambda)}=i_2(\lambda)
\]
ce qui exprime que les deux graduations \unsure{entières|partielles|naturelles|relatives} sont conjuguées complexes l'une de l'autre.

b) Il y a lieu de supposer que $\operatorname{Ob}C\ni Q(-1)$. Alors $Q(-1)$ \deleted{\unsure{est}} définit un $G$-module $F(Q(-1))$ de rang 1, d'où un hom
\[
\varepsilon:G\unsure{\longrightarrow|-}\mathrm G_{m\,K}
\]
Comme $Q(-1)$ est de \deleted{degré} poids 2, on aura
\[
\boxed{\varepsilon(\eta(\lambda))=\lambda^2}
\]

c) Supposons que sur le foncteur $F$ il y ait une filtration fonctorielle, compatible avec gr. Comment l'expliciter ? On considère
\[
F'=\mathrm{Gr}(F),
\]
% archive sheet number: 15

